\begin{figure}[htbp]
  \centering
  % Tự động thu phóng vừa với chiều ngang trang giấy
  \resizebox{\textwidth}{!}{
  \begin{tikzpicture}[
    % Hạ cỡ chữ xuống \footnotesize, tăng chiều rộng khối lên 2.8cm
    block/.style={rectangle, draw, thick, minimum width=2.8cm, minimum height=1.2cm, text width=2.6cm, align=center, fill=white, font=\footnotesize},
    decision/.style={diamond, draw, thick, minimum width=2.5cm, minimum height=2.5cm, text width=1.8cm, align=center, fill=white, font=\footnotesize},
    arrow/.style={draw, thick, -Stealth},
    ]

    % 1. VẼ KHUNG SWIMLANE (LÀN ĐƯỜNG)
    % Khung bao ngoài cùng (Mở rộng bao trọn toàn bộ các khối)
    \draw [thick] (-5.5, -2) rectangle (9, 10.5);
    
    % Đường kẻ dọc chia 3 làn (Mở rộng khoảng cách các làn)
    \draw [thick] (-1.5, -2) -- (-1.5, 10.5);
    \draw [thick] (4, -2) -- (4, 10.5);
    
    % Đường kẻ ngang chia phần Tiêu đề
    \draw [thick] (-5.5, 9.5) -- (9, 9.5);

    % Tiêu đề của 3 làn
    \node[font=\bfseries] at (-3.5, 10) {Người dân};
    \node[font=\bfseries] at (1.25, 10) {Hệ thống};
    \node[font=\bfseries] at (6.5, 10) {Đơn vị xử lý};

    % 2. ĐẶT CÁC KHỐI NODE THEO TỌA ĐỘ
    % Làn 1: Người dân (Tâm X = -3.5)
    \node [block] (phat_hien) at (-3.5, 8) {Phát hiện sự cố};
    \node [block] (phan_anh) at (-3.5, 6) {Phản ánh sự cố};

    % Làn 2: Hệ thống (Tâm X = 1.25)
    \node [block] (tiep_nhan) at (1.25, 8) {Tự động tiếp nhận sự cố};
    \node [block] (phan_loai) at (1.25, 6) {Phân loại mức độ ưu tiên};
    \node [block] (cap_nhat_td) at (1.25, 3) {Cập nhật tiến độ};
    \node [block] (luu_tru) at (1.25, -0.5) {Lưu trữ vào \gls{csdl}};

    % Làn 3: Đơn vị xử lý (Tâm X = 6.5)
    \node [block] (xac_thuc) at (6.5, 8) {Xác thực thông tin};
    \node [block] (thuc_hien) at (6.5, 6) {Thực hiện xử lý sự cố};
    \node [decision] (hoan_thanh) at (6.5, 3) {Hoàn thành};
    \node [block] (cap_nhat_tt) at (6.5, -0.5) {Cập nhật thông tin cơ sở hạ tầng};


    % 3. VẼ CÁC ĐƯỜNG MŨI TÊN
    % --- Nội bộ làn 1 ---
    \draw [arrow] (phat_hien) -- (phan_anh);

    % --- Làn 1 sang Làn 2 ---
    % Đi ngang sang phải (thêm 1cm) rồi gấp khúc đi thẳng lên khối Tiếp nhận
    \draw [arrow] (phan_anh.east) -- ++(1,0) |- (tiep_nhan.west);

    % --- Nội bộ làn 2 ---
    \draw [arrow] (tiep_nhan) -- (phan_loai);

    % --- Làn 2 sang Làn 3 ---
    \draw [arrow] (phan_loai.east) -- ++(1,0) |- (xac_thuc.west);

    % --- Nội bộ làn 3 ---
    \draw [arrow] (xac_thuc) -- (thuc_hien);
    \draw [arrow] (thuc_hien) -- (hoan_thanh);

    % --- Trả kết quả từ Quyết định ---
    % Trường hợp Sai: Từ Làn 3 chéo ngang về Làn 2
    \draw [arrow] (hoan_thanh.west) -- node[above,pos=0.6] {Sai} (cap_nhat_td.east);
    
    % Vòng lặp từ Cập nhật tiến độ (Làn 2) quay lại Thực hiện xử lý (Làn 3)
    % Cú pháp này đi lên trên 1 đoạn, rẽ ngang sang phải, rồi chui vào ĐÁY của khối "Thực hiện xử lý" để không đè lên mũi tên Phân loại
    \draw [arrow] (cap_nhat_td.north) -- ++(0, 1.2) -| ([xshift=-15pt]thuc_hien.south);

    % Trường hợp Đúng: Đi thẳng xuống
    \draw [arrow] (hoan_thanh.south) -- node[right] {Đúng} (cap_nhat_tt.north);

    % --- Từ Làn 3 về lại Làn 2 (Lưu trữ) ---
    \draw [arrow] (cap_nhat_tt.west) -- (luu_tru.east);

  \end{tikzpicture}
  }
  \caption{Quy trình tin học hóa Người dân phản ánh sự cố}
  \label{fig:quy_trinh_phan_anh}
\end{figure}